% block: computational-applications-heading page:43 type:chapter
\clearpage
\section{{Some Applications to Computational Science and Data Science}}\label{chap:computational-applications}
\setcounter{footnote}{0}
% block: computational-applications-abstract page:43 type:ordinary-paragraph
\noindent\textbf{{Abstract}}\quad
{Linear systems of the form $A\boldsymbol{x}=\boldsymbol{b}$ arise in a rich variety of applications such as computational science and data science. For computational science, physical phenomena are often described as partial differential equations. For solving partial differential equations, the finite-difference methods and the finite element method are widely used, leading to a problem of solving linear systems. Large-scale electronic structure calculation for condensed matter physics and lattice quantum chromodynamics for particle physics require solving large-scale linear systems. For data science, regression and classification are fundamental tasks that also require solving linear systems. Minimizing or maximizing functions is one of the most important optimization problems, which requires solving linear systems when Newton's method is used. We will see in this chapter how the linear systems arise in these applications by using simple and concrete examples.}
